\documentclass[10pt,a4paper]{amsart}
\usepackage[margin=19mm]{geometry}
\usepackage{amsmath,amssymb,mathtools}
\usepackage[scr=rsfs,cal=euler]{mathalfa}
\usepackage{xfrac}
\usepackage[unicode,hidelinks,bookmarks=false]{hyperref}
\newcommand{\ie}{i.e.\ }
\newcommand{\cf}{cf.\ }
\newcommand{\resp}{respectively\ }
\newcommand{\Subsection}{\subsection}
\providecommand{\nobreakdash}{\nobreak\hbox{-}}
\setlength{\emergencystretch}{2em}
\multlinegap0pt
\usepackage{amssymb}
\usepackage{enumerate}
\usepackage[shortlabels]{enumitem}
\usepackage{tikz}
\usepackage{cancel}
\newtheorem{cor}{Corollary}
\newtheorem{lem}{Lemma}
\newtheorem{prop}{Proposition}
\newcommand{\conv}[1]{\star_{#1}}
\title{On the Geometry and Optimization of Polynomial Convolutional Networks}
\author{Vahid Shahverdi, Giovanni Luca Marchetti, Kathl\'en Kohn}
\date{Source: arXiv:2410.00722v3 (CC BY 4.0)}
\begin{document}
\maketitle

\noindent\emph{Source and licence.} On the Geometry and Optimization of Polynomial Convolutional Networks. Version arXiv:2410.00722v3 (CC BY 4.0), distributed by arXiv under the Creative Commons Attribution 4.0 International licence, http://creativecommons.org/licenses/by/4.0/, as recorded in the arXiv licence metadata for this exact version and shown on the abstract page https://arxiv.org/abs/2410.00722v3. Full licence notice: \texttt{licenses/arxiv-2410.00722-CC-BY-4.0.txt}.\par\smallskip

\noindent\textit{Editorial scope.} Counted unit: the article's prop prop:kerdiff (source0.tex lines 663--665). The article's complete original proof of it is delivered with it (source0.tex lines 666--701, one \texttt{proof} environment of 36 lines). No mathematics is written, repaired or re-derived: the statement and the proof are the article's own lines, in the article's own notation, and no retained line is substituted or re-flowed.  Retained uncounted as the source-local context the proof needs: lines 326--328; lines 653--658.  Nothing was omitted from any retained span, so the package carries no omission record.  No retained line contains a citation, so the wrapper carries no bibliography and no bibliography entry.  No retained line contains a figure environment, an image-inclusion command or drawing code, so no image is delivered and the package declares no asset dependency.  Editorial formatting only, in addition: an AMS \texttt{amsart} preamble that restates the article's own declarations and macros used by the retained lines, namely the environments \texttt{cor}, \texttt{lem}, \texttt{prop} on separate counters, together with the standard packages those lines need.  The retained environments print in this document as Corollary 1, Lemma 1, Proposition 1 in order of appearance, whereas the article prints the same items with the numbering of its own full document.  The complete native source of the article as distributed in the arXiv e-print for this exact version is bundled unchanged as \texttt{sources/166101/source0.tex}.
\begin{cor}\label{cor:basepts}
    If $\varphi_\mathbf{w}(x) = 0$ for all $x \in \mathbb{C}^{d_0}$, then $w_i = 0$ for some $i$. 
\end{cor}

\begin{lem}\label{lemm:homogendiff}
Given $\lambda_0, \ldots, \lambda_{L-1} \in \mathbb{C}$, the differential $\textnormal{J}_\mathbf{w} \varphi$ sends $(\lambda_{0}w_{0}, \ldots, \lambda_{L-1} w_{L-1})$ to:
\begin{equation}
    (\lambda_{L-1} + r \lambda_{L-2} + \cdots + r^{L-1}\lambda_0) \ \varphi_\mathbf{w}. 
\end{equation}
\end{lem}

\begin{prop}\label{prop:kerdiff}
Suppose that $w_i \not = 0 $ for all $i$. Then the kernel of $\textnormal{J}_\mathbf{w} \varphi$ is generated, as a linear subspace of $\mathbb{C}^{|\mathbf{k}|}$, by $(0,\ldots,   0, w_{L-2}, -rw_{L-1})$, $(0, \ldots, 0, w_{L-3}, -rw_{L-2},0)$, $\ldots$, $(w_0, -rw_1, 0, \ldots, 0)$. 
\end{prop}

\begin{proof}
We proceed by induction over the number of layers $L$. For $L=1$,  $\varphi_\mathbf{w}(x) = w_0 \conv{s_0} x$, which is linear in $w_0$ and does not vanish for all $x$ since $w_0 \not = 0$. The differential is therefore injective, and its kernel is trivial. 

Suppose now that $L > 1$. In order to compute the kernel $\textnormal{J}_\mathbf{w} \varphi$, we need to solve
\begin{equation}\label{eq:kerndiff}
   \dot{w}_{L-1} \conv{s_{L-1}} \sigma_r\left(  \varphi_{\mathbf{w}'} \right) + r w_{L-1} \conv{s_{L-1}} \sigma_{r-1}\left(  \varphi_{\mathbf{w}'} \right) \odot \left( \textnormal{J}_\mathbf{w'} \varphi \right)(\dot{\mathbf{w}}') = 0. 
\end{equation}
If $\dot{\mathbf{w}}'$ belongs to the kernel of $\textnormal{J}_\mathbf{w'} \varphi$, then since $\varphi_{\mathbf{w}'} \not = 0$ by hypothesis (see Corollary \ref{cor:basepts}), we must have $\dot{w}_{L-1} = 0$. It follows from the inductive hypothesis that $\dot{w}$ is a linear combination of the tuples of filters in the statement. We therefore assume that $\left( \textnormal{J}_\mathbf{w'} \varphi\right)(\dot{\mathbf{w}}') \not = 0$. Moreover, without loss of generality, we assume that  $w_{L-1}[0] \not = 0$. 

We will now decompose $\varphi_\mathbf{w}$ as a sum of CNNs with smaller filter size on the last layer. To this end, consider the tuple of filters:
\begin{equation}
\mathbf{w}^+=(w_{0}, \ldots, w_{L-2}, (w_{L-1}[0])), \hspace{4em} \mathbf{w}^-=(w_{0}, \ldots, w_{L-2}, (0, w_{L-1}[1], \ldots, w_{L-1}[k_{L-1} -1])).
\end{equation}
Then $\varphi_\mathbf{w} = \varphi_{\mathbf{w}^+} + \varphi_{\mathbf{w}^-}$. Now, let $i$ be the minimal index such that $x[i]$ appears in $\varphi_{\mathbf{w}'}[0]$, seen as a scalar-valued polynomial in the input variable $x$. Recall the basic equivariance property of CNNs, meaning that shifting indices in the output variable is equivalent to shifting them in the input one (assuming the latter shift accounts for the strides). This implies that $x[i]$ does not appear in neither $\left( \textnormal{J}_\mathbf{w^-} \varphi \right)(\dot{\mathbf{w}}^-)$ nor $\varphi_{\mathbf{w}^-}$. Since $0 = \textnormal{J}_\mathbf{w} \varphi = \textnormal{J}_{\mathbf{w}^+} \varphi + \textnormal{J}_{\mathbf{w}^-} \varphi$, $x[i]$ does not appear in $\textnormal{J}_{\mathbf{w}^+} \varphi$ as well. Write: 
\begin{equation}
\varphi_{\mathbf{w}'}[0] = x[i]f + g, \hspace{7em} \left( \textnormal{J}_\mathbf{w'} \varphi \right)(\dot{\mathbf{w}}')[0] = x[i]a + b,
\end{equation}
where $f,g,a,b$ are polynomials such that $f \not =  0$ and $x[i]$ does not appear in $g$ nor $b$. The (vanishing) terms containing $x[i]$ in $\textnormal{J}_{\mathbf{w}^+} \varphi$ have the form:
\begin{equation}
    0 = \dot{w}_{L-1}[0] \left((x[i]f + g)^r - g^r \right) + rw_{L-1}[0]\left((x[i]f + g)^{r-1}(x[i]a + b) - g^{r-1}b\right).
\end{equation}
By manipulating the above expression, we obtain:
\begin{equation}
(x[i]f + g)^{r-1}\left(\dot{w}_{L-1}[0](x[i]f + g) + rw_{L-1}[0](x[i]a + b) \right) = g^{r-1}( \dot{w}_{L-1}[0]g + r w_{L-1}[0]b).
\end{equation}
The right-hand side of the above equation does not contain $x[i]$. Since $f \not = 0$, it follows that $\dot{w}_{L-1}[0](x[i]f + g) + rw_{L-1}[0](x[i]a + b) = 0$, implying:
\begin{equation}
\varphi_{\mathbf{w}'}[0] = x[i]a + b = \lambda \ (x[i]f + g) = \lambda \left( \textnormal{J}_\mathbf{w'} \varphi \right)(\dot{\mathbf{w}}')[0], \hspace{3em} \lambda := -\frac{\dot{w}_{L-1}[0]}{w_{L-1}[0]}.
\end{equation}
By looking at (the $0$-th entry of) Equation \ref{eq:kerndiff}, we deduce $\dot{w}_{L-1} = -r \lambda w_{L-1}$. By Lemma \ref{lemm:homogendiff}, $\dot{\mathbf{w}}' = (\lambda_{0}w_{0}, \ldots, \lambda_{L-2}w_{L-2})$ for some $\lambda_i \in \mathbb{C} \setminus \{0\}$ such that $\lambda_{L-2} + r \lambda_{r-3} + \cdots + r^{L-2} \lambda_0 = \lambda$. But then
\begin{align}
\dot{\mathbf{w}} &= (\lambda_0 w_0, \ldots,  \lambda_{L-2}w_{L-2}, -r \lambda w_{L-1}) = \\
&= \lambda (0, \ldots, 0, w_{L-2}, -rw_{L-1}) + (\lambda_{0} w_0, \ldots, \lambda_{L-3}w_{L-3}, (\lambda_{L-2} - \lambda)w_{L-2}, 0).
\end{align}
Since $(\lambda_{L-2} - \lambda) + r \lambda_{L-3} + \cdots + r^{L-2}\lambda_0 = 0$, it follows from Lemma \ref{lemm:homogendiff} that the second summand of the above expression belongs to the kernel of $ \textnormal{J}_\mathbf{w'} \varphi$, and the desired result follows from the inductive hypothesis.
\end{proof}

\end{document}
